% TOP_PROBLEM: {"id": 1100306, "title": "Questions in Geometric Group Theory — Q 3.6", "category_id": 17, "category": {"id": 17, "name": "group_theory", "display_name": "Group Theory", "description": "Problems about groups, group actions, representations, and related algebraic structures.", "slug": "group-theory", "order_index": 17, "created_at": "2026-07-31T15:26:25.671Z"}, "status": "partially_solved", "problem_number": "AMR-010-0306", "set_id": 13}
% FIRST_POSTED: 2026-10-01
% ENTRY_KIND: statement_only
% UnsolvedMath ID 1100306; source catalogue code AMR-010-0306
% !TeX program = lualatex
% Definition and sources only; this file is not an LLM solution attempt.
\documentclass[11pt,a4paper]{article}
\usepackage[margin=25mm]{geometry}
\usepackage{fontspec}
\setmainfont{lmroman10-regular.otf}[BoldFont=lmroman10-bold.otf,
 ItalicFont=lmroman10-italic.otf,BoldItalicFont=lmroman10-bolditalic.otf]
\usepackage{amsmath,amssymb,mathtools}
\usepackage{unicode-math}
\setmathfont{latinmodern-math.otf}
\AtBeginDocument{\let\setminus\smallsetminus}
\usepackage{xurl,hyperref}
\hypersetup{colorlinks=true,urlcolor=blue}
\setcounter{secnumdepth}{0}
\setlength{\parindent}{0pt}
\setlength{\parskip}{5pt}
\setlength{\emergencystretch}{3em}
\begin{document}
\section*{Questions in Geometric Group Theory — Q 3.6}\label{1100306}
{\small UnsolvedMath ID 1100306 · AMR-010-0306 · Group Theory · AMR Open Problem Lists}
\subsection{Definitions and mathematical statement}
Let $F_m$ and $F_n$ be free groups of finite positive ranks, on separate bases. Choose nonidentity elements $u\in F_m$ and $v\in F_n$. Identify the generators of an infinite cyclic group with $u$ and $v$, respectively, and form the amalgamated free product
\[
 G=F_m*_{\mathbb Z}F_n
   =\langle x_1,\ldots,x_m,y_1,\ldots,y_n\mid u(x)=v(y)\rangle.
\]
The nonidentity assumptions ensure that both cyclic edge maps are injective. Either word may be a proper power, and no maximal-cyclic condition is imposed.

A finitely generated group is a limit group here if it is fully residually free: for every finite subset $E$ there is a homomorphism to a nonabelian free group whose restriction to $E$ is injective. The map may depend on $E$. Equivalently, every finite set of nonidentity elements can simultaneously survive under one homomorphism to a free group.

Characterize exactly which choices of $m,n,u,v$ produce a limit group. A complete answer should give necessary and sufficient conditions on the amalgamating words, allowing changes of free bases and other choices that give isomorphic amalgams. Ordinary residual freeness, which separates one element at a time, is not the definition being asked for. The desired characterization includes proper-power attachments as well as root-free words. It must determine the full group obtained from the amalgam, rather than only whether the two displayed free factors embed.
\subsection{Short English statement}
Which cyclic amalgams of two finite-rank free groups are fully residually free?
\subsection{Sources}
\begin{itemize}
\item[S1] ULAM.AI and the UnsolvedMath Contributors, supplied problems.json, record 1100306 (AMR-010-0306), AMR Open Problem Lists. Catalogue identity and source transcription; CC BY 4.0.
\url{https://huggingface.co/datasets/ulamai/UnsolvedMath}.
\item[S2] Bestvina - Questions in Geometric Group Theory (pdf) (2004), Question 3.6 (PDF page 11). Original source indexed by AMR-010-0306.
\url{https://www.math.utah.edu/~bestvina/eprints/questions-updated.pdf}.
\end{itemize}
\end{document}
